\begin{textsample}{POS Dim 1 – go – Score 36.00 – go/go\_em\_50.txt}  \label{ex:f1_pos_109}
3.1 - Ciências da Natureza e Suas Tecnologias e o DC-GOEM As Ciências da Natureza têm \textbf{desempenhado} \textbf{um} papel cada vez mais importante na vida em sociedade.
As bases científicas \textbf{historicamente} lançadas sustentam muitos conhecimentos \textbf{que} \textbf{hoje} tomamos como certos e \textbf{que} têm \textbf{influenciado} diretamente nos avanços técnico-científicos.
Atualmente, os impactos positivos gerados pelos estudos nessa área podem ser percebidos, por exemplo, nos medicamentos \textbf{que} salvam vidas, nas tecnologias inseridas nos lugares em \textbf{que} \textbf{vivemos} e trabalhamos, na forma como nos comunicamos e numa infinidade de outras aplicações \textbf{que} irão moldar, inclusive, a maneira como interpretamos o mundo. \textbf{Contudo}, apesar de todas essas evidências sobre a importância do conhecimento e desenvolvimento científico, ainda pode -se perguntar : “ Por \textbf{que} estudar Ciências da Natureza no Ensino Médio? ” Para responder a esse questionamento, é importante \textbf{que} se entenda \textbf{um} dos objetivos da Ciências da Natureza e suas Tecnologias : a compreensão do mundo natural e a avaliação dos impactos ambientais gerados pelas ações antrópicas, a partir de \textbf{um} olhar articulado, entre a Biologia, a Química e a Física.
Nesse sentido, os conhecimentos da área de Ciências da Natureza e suas Tecnologias podem ser mobilizados para desenvolver nos/nas estudantes a capacidade de interpretar o mundo \textbf{que} os/as cerca, tornando-os/as \textbf{sujeitos} pensantes e críticos, para \textbf{que} se formem cidadãos/cidadãs implicados/as no desenvolvimento social, político e econômico do país.
Voltando nosso olhar para o passado, podemos concluir \textbf{que}, desde os primórdios, a \textbf{humanidade} busca compreender a natureza por meio da manipulação e domínio de fenômenos naturais para o aperfeiçoamento das técnicas de sobrevivência.
Assim, a Ciência se desenvolveu junto e por causa da espécie humana e percorreu \textbf{um} longo caminho até a formação de como a conhecemos \textbf{hoje} e ainda está em constante evolução.
O mundo passou por \textbf{inúmeras} transformações, desde os filósofos gregos pautados no empirismo, até o desenvolvimento dos procedimentos científicos \textbf{que} trouxeram a \textbf{humanidade} à sociedade tecnológica \textbf{atual}.
Tais transformações, impulsionadas a partir do \textbf{século} VII, tiveram como fruto mais recente a ciência \textbf{moderna} \textbf{focada} na experimentação como forma de conhecer e delimitar leis naturais.
Nascia, então, essa ciência \textbf{moderna}, apoiada em grandes \textbf{mentes} como as de Lavoisier e Galileu Galilei.
Eles \textbf{argumentavam} \textbf{que} somente por meio de experimentos quantitativos poderíamos testar hipóteses e, seguramente, \textbf{desvendar} as lógicas do mundo ( PORTO, 2008 ).
Nesse contexto, concluímos \textbf{que} a Ciência, de modo geral, pode ser \textbf{uma} atividade dinâmica ou metódica e ainda criativa com \textbf{uma} longa e \textbf{interessante} história, em \textbf{que} várias sociedades, em diferentes épocas, contribuíram para a construção do conhecimento, além do desenvolvimento e entendimento científico. \textbf{Cientistas} continuamente avaliam a solidez do conhecimento científico, testando leis e \textbf{teorias} anteriormente estabelecidas, modificando -as à medida \textbf{que} aparecem novas evidências ou mesmo ressignificando -as, a partir de evidências já existentes. ( SINGH, 2006 ).
A ciência racionaliza o processo de pesquisa \textbf{que} se inicia por meio de indagações direcionadoras, sendo estas, muitas vezes, mais importantes \textbf{que} as próprias respostas a serem obtidas.
O ato de questionar é de \textbf{suma} importância no processo de construção da ciência, mas esta não sobrevive nem se dissemina sem \textbf{que} seja ensinada ( SILVA ; FERREIRA ; VIEIRA 2017 ).
Assim, no processo de apresentação e \textbf{contextualização} dos conhecimentos da área de Ciências da Natureza e suas Tecnologias, a escola ocupa \textbf{um} lugar \textbf{central}.
O sucesso dessa \textbf{tarefa} pode estar em aproveitar \textbf{uma} característica \textbf{que} os seres humanos possuem : a curiosidade!
Ao longo da Educação Básica, a curiosidade pode ser \textbf{instigada} e trabalhada juntamente com as habilidades socioemocionais.
A partir da intencionalidade pedagógica do/a professor/a, habilidades mais complexas e híbridas, como a criatividade e o pensamento crítico, podem ser, inicialmente, trabalhadas no Ensino Fundamental e continuamente desenvolvidas no Ensino Médio, favorecendo \textbf{uma} melhor articulação entre essas etapas e evitando -se rupturas.
Segundo Moreira ( 2011 ), os/as estudantes trazem consigo interesses individuais, experiências pessoais e culturais diversas, impactando o seu conhecimento prévio sobre ciência, tecnologia e o mundo em \textbf{que} \textbf{vivem}.
Tal experiência e conhecimento podem ser alargados quando o professor estimula questionamentos e a construção conjunta do conhecimento.
A partir do acolhimento e da escuta ativa dos/as estudantes, tanto as habilidades cognitivas quanto as habilidades socioemocionais podem ser desenvolvidas, promovendo o protagonismo dos/as estudantes e fazendo-os/as reconhecer como participantes ativos/as das sociedades as quais estão inseridos/as.
Em busca da formação integral do ser humano, o DC-GOEM - área Ciências da Natureza e suas Tecnologias - traz, em sua essência, objetivos de aprendizagem específicos \textbf{que} foram construídos, estruturados e organizados tendo como referência a Base Nacional Comum Curricular do Ensino Médio.
O objetivo destes documentos é garantir meios para o desenvolvimento das competências e habilidades em \textbf{um} contexto social, ético, crítico e científico global.
O DC-GOEM oferece meios de favorecer e possibilitar o planejamento das aulas por parte dos/as professores/as, garantindo o desenvolvimento das habilidades \textbf{que} estão ligadas às três competências específicas de Ciências da Natureza e suas Tecnologias para o Ensino Médio, as quais, por sua vez, relacionam -se diretamente com as dez competências gerais da BNCC.
Os meios em questão são intencionalidades didáticas na forma de objetivos de aprendizagem.
Estes podem subsidiar o trabalho dos/as professores/as apresentando os conhecimentos escolares essenciais \textbf{que} devem ser trabalhados em sala de aula, articulando -os a procedimentos didáticos e metodologias, além de sugerirem finalidades, sejam estas para a própria área, áreas afins ou para a vida.
Descrevem conceitos, conhecimentos e processos, ou seja, a aprendizagem esperada dos/as estudantes em cada habilidade a ser desenvolvida.
Isso oferece \textbf{um} quadro de progressão e desenvolvimento dentro da área de Ciências da Natureza e suas Tecnologias, o \textbf{que} ajuda a planejar e monitorar a aprendizagem, e fazer análises sobre o desempenho dos/as estudantes ( ANDERSON \textbf{et} al., 2001 ).
As habilidades específicas da área de Ciências da Natureza e suas Tecnologias fornecidas pela BNCC do Ensino Médio foram \textbf{aqui} reescritas na forma de objetivos de aprendizagem, tendo como recurso de formulação a taxonomia de Bloom ( ANDERSON \textbf{et} al., 2001 ).
Esses objetivos constituem em \textbf{uma} evolução na busca de \textbf{uma} alternativa ao modelo \textbf{embasado} na aprendizagem de saberes disciplinares organizados ao redor de matérias convencionais, na qual o/a estudante deveria assumir os conteúdos como eram definidos pelas diferentes propostas científicas ( ZABALA ; ARNAU, 2014 ).
Considerando -se a interação entre os componentes curriculares : Biologia, Física e Química, os objetivos de aprendizagem estão organizados em três partes \textbf{que} evidenciam as habilidades cognitivas a serem aplicadas.
A primeira parte faz referência a “ o quê ” se deseja aprender.
A seguinte trata do “ como ”, quer \textbf{dizer}, a metodologia \textbf{que} pode ser empregada.
A última refere -se ao “ para quê ”, ou seja, a finalidade do “ o quê ” e do “ como ” aprender.
Cada objetivo de aprendizagem segue o comando implícito “ os/as estudantes irão... ”, apresentando, inicialmente, \textbf{um} verbo no infinitivo, \textbf{que} indica o processo cognitivo a ser desenvolvido, seguido do conhecimento \textbf{que} se espera \textbf{que} o/a estudante mobilize para alcançar o objetivo.
Em seguida, \textbf{um} verbo no gerúndio especifica padrões, condições ou critérios de desempenho esperados, esclarecendo mais detalhadamente sobre o processo de aprendizagem.
E, por fim, mais \textbf{um} verbo no infinitivo, \textbf{que} traz \textbf{uma} justificativa a todo o processo e abre possibilidades de ampliação do conhecimento do/a estudante.
Esses três verbos obedecem a \textbf{uma} determinada gradação ascendente de processos cognitivos fundamentais \textbf{que} podem ser agrupados nas \textbf{categorias} de nível básico, operacional ou global, expressando \textbf{uma} evolução na complexidade de habilidades de pensamento.
Dessa forma, para \textbf{que} o processo de \textbf{ensino-aprendizagem} aconteça, as habilidades e competências a serem alcançadas pelos/as estudantes deverão ir além da \textbf{memorização} de fatos isolados ou da repetição de termos e \textbf{teorias} específicas.
Por isso, é necessária a construção de objetivos de aprendizagem relacionados ao pensamento crítico-científico, de modo a encontrar \textbf{aplicabilidade} tanto no contexto científico, quanto no cotidiano do/a estudante, além de oferecer a base necessária para continuar seus estudos posteriores.
Espera -se \textbf{que} os/as estudantes se sintam mais envolvidos/as quando são capazes de perceber a conexão entre os conceitos científicos \textbf{que} aprenderam e a aplicação no mundo a sua volta.
Todos esses aspectos do DC-GOEM da área de Ciências da Natureza e suas Tecnologias têm como objetivo auxiliar os/as profissionais da educação de nosso estado na formação e desenvolvimento intencional de vários aspectos dos/as estudantes, como cognitivos e socioemocionais, para \textbf{que} todos/as sejam capazes de agir de maneira autônoma, como \textbf{sujeitos} e cidadãos/ãs protagonistas em seus projetos de vida.

% matched lemmas: aplicabilidade, aqui, argumentar, atual, categoria, central, cientista, contextualização, contudo, desempenhar, desvendar, dizer, embasar, ensino-aprendizagem, et, focar, historicamente, hoje, humanidade, influenciar, instigar, interessante, inúmero, memorização, mente, moderno, que, sujeito, sumo, século, tarefa, teoria, um, vivar, viver
\end{textsample}
